\documentclass{article}

% packages
\usepackage{amsmath, amsthm, thmtools, amsfonts, amssymb, luacode, catchfile, tikzducks, hyperref, ifthen}
\ifcsname c@kobocompile\endcsname
	\usepackage[a5paper, total={1072pt, 1448pt}, margin=10pt, includeheadfoot]{geometry} % set page margins
\else
	\usepackage[a4paper, margin=50pt, includeheadfoot]{geometry}
\fi
\usepackage[shortlabels]{enumitem}
\usepackage[skip=3pt, indent=0pt]{parskip}

% language
\usepackage[bidi=basic, layout=tabular, provide=*]{babel}
\ifcsname c@english\endcsname
	\babelprovide[main, import]{english}
\else
	\babelprovide[main, import]{hebrew}
	\babelprovide{rl}
\fi
%\babelfont{rm}{Libertinus Serif}
\babelfont{rm}[Renderer=Harfbuzz]{Libertinus Serif}
\babelfont{sf}{Libertinus Sans}
\babelfont{tt}{Libertinus Mono}

% style
\AddToHook{cmd/section/before}{\clearpage}	% Add line break before section
\linespread{1.3}
\setcounter{secnumdepth}{0}		% Remove default number tags from sections, this won't do well with theorems
\AtBeginDocument{\setlength{\belowdisplayskip}{3pt}}
\AtBeginDocument{\setlength{\abovedisplayskip}{3pt}}
\graphicspath{ {../images/} }

% operators
\DeclareMathOperator\cis{cis}
\DeclareMathOperator\Sp{Sp}
\DeclareMathOperator\tr{tr}
\DeclareMathOperator\im{Im}
\DeclareMathOperator\re{Re}
\DeclareMathOperator\diag{diag}
\DeclareMathOperator*\lowlim{\underline{lim}}
\DeclareMathOperator*\uplim{\overline{lim}}
\DeclareMathOperator\rng{rng}
\DeclareMathOperator\Sym{Sym}
\DeclareMathOperator\Arg{Arg}
\DeclareMathOperator\Log{Log}
\DeclareMathOperator\dom{dom}
\DeclareMathOperator\supp{Supp}
\DeclareMathOperator\var{Var}
\DeclareMathOperator\cov{Cov}

% commands
%\renewcommand\qedsymbol{\textbf{מש''ל}}
%\renewcommand\qedsymbol{\fbox{\emoji{lizard}}}
\newcommand{\Aa}[0]{\mathcal{A}}
\newcommand{\Bb}[0]{\mathcal{B}}
\newcommand{\CC}[0]{\mathbb{C}}
\newcommand{\Cc}[0]{\mathcal{C}}
\newcommand{\EE}[0]{\mathbb{E}}
\newcommand{\FF}[0]{\mathbb{F}}
\newcommand{\Ff}[0]{\mathcal{F}}
\newcommand{\Ii}[0]{\mathcal{I}}
\newcommand{\Gg}[0]{\mathcal{G}}
\newcommand{\Ll}[0]{\mathcal{L}}
\newcommand{\Mm}[0]{\mathcal{M}}
\newcommand{\NN}[0]{\mathbb{N}}
\newcommand{\Nn}[0]{\mathcal{N}}
\newcommand{\PP}[0]{\mathbb{P}}
\newcommand{\Pp}[0]{\mathcal{P}}
\newcommand{\QQ}[0]{\mathbb{Q}}
\newcommand{\RR}[0]{\mathbb{R}}
\newcommand{\Rr}[0]{\mathcal{R}}
\newcommand{\Ss}[0]{\mathcal{S}}
\newcommand{\TT}[0]{\mathbb{T}}
\newcommand{\Uu}[0]{\mathcal{U}}
\newcommand{\Vv}[0]{\mathcal{V}}
\newcommand{\Ww}[0]{\mathcal{W}}
\newcommand{\ZZ}[0]{\mathbb{Z}}
\newcommand{\acts}[0]{\circlearrowright}
\newcommand{\explain}[2] {
	\begin{flalign*}
		 && \text{#2} && \text{#1}
	\end{flalign*}
}
\newcommand{\maketitleprint}[0]{ \begin{center}
	%\begin{tikzpicture}[scale=3]
	%	\duck[graduate=gray!20!black, tassel=red!70!black]
	%\end{tikzpicture}	
	\includegraphics[width=6cm]{cover}
\end{center}
}

% theorem commands
\newtheoremstyle{c_remark}
	{}	% Space above
	{}	% Space below
	{}% Body font
	{}	% Indent amount
	{\bfseries}	% Theorem head font
	{}	% Punctuation after theorem head
	{.5em}	% Space after theorem head
	{\thmname{#1}\thmnumber{ #2}\thmnote{ \normalfont{\text{(#3)}}}}	% head content
\newtheoremstyle{c_definition}
	{3pt}	% Space above
	{3pt}	% Space below
	{}% Body font
	{}	% Indent amount
	{\bfseries}	% Theorem head font
	{}	% Punctuation after theorem head
	{.5em}	% Space after theorem head
	{\thmname{#1}\thmnumber{ #2}\thmnote{ \normalfont{\text{(#3)}}}}	% head content
\newtheoremstyle{c_plain}
	{3pt}	% Space above
	{3pt}	% Space below
	{\itshape}% Body font
	{}	% Indent amount
	{\bfseries}	% Theorem head font
	{}	% Punctuation after theorem head
	{.5em}	% Space after theorem head
	{\thmname{#1}\thmnumber{ #2}\thmnote{ \text{(#3)}}}	% head content

\ifcsname c@english\endcsname
	\theoremstyle{plain}
	\newtheorem{theorem}{Theorem}[section]
	\newtheorem{lemma}[theorem]{Lemma}
	\newtheorem{proposition}[theorem]{Proposition}
	\newtheorem*{proposition*}{Proposition}
	%\newtheorem{corollary}[theorem]{אין חלופה עברית}

	\theoremstyle{definition}
	\newtheorem{definition}[theorem]{Definition}
	\newtheorem*{definition*}{Definition}
	\newtheorem{example}{Example}[section]
	\newtheorem{exercise}{Exercise}[section]

	\theoremstyle{remark}
	\newtheorem*{remark}{Remark}
	\newtheorem*{solution}{Solution}
	\newtheorem{conclusion}[theorem]{Conclusion}
	\newtheorem{notation}[theorem]{Notation}
\else
	\theoremstyle{c_plain}
	\newtheorem{theorem}{משפט}[section]
	\newtheorem{lemma}[theorem]{למה}
	\newtheorem{proposition}[theorem]{טענה}
	\newtheorem*{proposition*}{טענה}
	%\newtheorem{corollary}[theorem]{אין חלופה עברית}

	\theoremstyle{c_definition}
	\newtheorem{definition}[theorem]{הגדרה}
	\newtheorem*{definition*}{הגדרה}
	\newtheorem{example}{דוגמה}[section]
	\newtheorem{exercise}{תרגיל}[section]

	\theoremstyle{c_remark}
	\newtheorem*{remark}{הערה}
	\newtheorem*{solution}{פתרון}
	\newtheorem{conclusion}[theorem]{מסקנה}
	\newtheorem{notation}[theorem]{סימון}
\fi

% Questions related commands
\newcounter{question}
\setcounter{question}{1}
\newcounter{sub_question}
\setcounter{sub_question}{1}

\ifcsname c@english\endcsname
	\newcommand{\question}[1][0]{
		\ifthenelse{#1 = 0}{}{\setcounter{question}{#1}}
		\section{Question \arabic{question}}
		\addtocounter{question}{1}
		\setcounter{sub_question}{1}
	}

	\newcommand{\subquestion}[1][0]{
		\ifthenelse{#1 = 0}{}{\setcounter{sub_question}{#1}}
		\subsection{Part \alph{sub_question}}
		\addtocounter{sub_question}{1}
	}
\else
	\newcommand{\question}[1][0]{
		\ifthenelse{#1 = 0}{}{\setcounter{question}{#1}}
		\section{שאלה \arabic{question}}
		\addtocounter{question}{1}
		\setcounter{sub_question}{1}
	}

	\newcommand{\subquestion}[1][0]{
		\ifthenelse{#1 = 0}{}{\setcounter{sub_question}{#1}}
		\subsection{סעיף \localecounter{letters.gershayim}{sub_question}}
		\addtocounter{sub_question}{1}
	}
\fi

% import lua and start of document
\directlua{common = require ('../common')}

\GetEnv{AUTHOR}

% headers
\author{\AUTHOR}
\date\today

\title{פתרון מטלה 10 --- אנליזה פונקציונלית, 80417}

\DeclareMathOperator\sgn{sgn}

\begin{document}
\maketitle
\maketitleprint[teal]

\question{}
לכל $f, g : \RR \to \RR$ פונקציות אינטגרביליות ו־$2 \pi$־מחזוריות נגדיר,
\[
	(f * g)(x)
	= \int_{-\pi}^{\pi} f(x - u) g(u)\ du
\]

\subquestion{}
נראה שאם $g \in \tilde{C}[-\pi, \pi]$ אז $f * g \in C[-\pi, \pi]$.
\begin{proof}
	מהנתון $g$ רציפה, $2 \pi$־מחזורית ובפרט אינטגרבילית רימן.
	עוד נתון כי $f$ אינטגרבילית, לאו דווקה רציפה, אך מחזורית.
	נבחין כי $f(x - u) g(u)$ היא פונקציה אינטגרבילית ומתקיים,
	\[
		(f * g)(x)
		= \int_{-\pi}^{\pi} f(x - u) g(u)\ du
	\]
	היא פונקציה רציפה בכל נקודה, ואילו היא לא רציפה בנקודה, אז זוהי נקודת אי־רציפות מסדר שני ובפרט גם ל־$|f|$ בהכרח שואפת לאינסוף באחד מהקצוות של הנקודה הזו.
	אבל במקרה זה היא לא אינטגרבילית רימן.
\end{proof}

\subquestion{}
נראה שאם $g \in \tilde{C}^1[-\pi, \pi]$ אז $f * g \in C^1[-\pi, \pi]$,
וכן שלכל $x \in [-\pi, \pi]$ מתקיים,
\[
	(f * g)'(x)
	= (f * g')(x)
\]
\begin{proof}
	אם $F(x) = \int_{-\pi}^{x} f(u)\ du$,
	\[
		\int_{-\pi}^{\pi} f(x - u) g(u)\ du
		= \left. -F(x - u) g(u) \right\lvert_{u = -\pi}^{u = \pi} - \int_{-\pi}^{\pi} -F(x - u) g'(u)\ du
		= 0 + \int_{-\pi}^{\pi} F(x - u) g'(u)\ du
	\]
	כלומר,
	\[
		(f * g)(x)
		= \int_{-\pi}^{\pi} F(x - u) g'(u)\ du
		= (F * g')(x)
		= (g' * F)(x)
	\]
	$F$ רציפה כקדומה ו־$g'$ רציפה שכן $g \in C^1[-\pi, \pi]$.
	לכן,
	\begin{align*}
		(f * g)'(x)
		& = \lim_{h \to 0} \frac{1}{h} ((f * g)(x + h) - (f * g)(x)) \\
		& = \lim_{h \to 0} \frac{1}{h} \int_{-\pi}^{\pi} f(x - u) g(u) - f(x - u - h) g(u + h)\ du \\
		& = \lim_{h \to 0} \frac{1}{h} \int_{-\pi}^{\pi} f(x - u) g(u) - f(x - u - h) g(u) + f(x - u - h)( g(u) - g(u + h))\ du \\
		& = \lim_{h \to 0} \frac{1}{h} \int_{-\pi}^{\pi} (f(x - u) - f(x - u - h)) g(u) + f(x - u - h)( g(u) - g(u + h))\ du \\
		& = \lim_{h \to 0} \frac{1}{h} ((f(x - u) - f(x - u - h)) * g(u)) + \int_{-\pi}^{\pi} f(x - u) g'(u)\ du \\
		& = \lim_{h \to 0} \frac{1}{h} ((F(x - u) - F(x - u - h)) * g'(u)) + (f * g')(x) \\
		& = (f * g')(x)
	\end{align*}
	והטענה נובעת ממחזוריות ומגזירות הפונקציה הקדומה של $f$.
\end{proof}

\question{}
\subquestion{}
תהי ${\{ x_n \}}_{n = 1}^\infty \subseteq \RR$ כך שלכל $n$ מתקיים,
\[
	|x_{n + 1} - x_n| \le \frac{C}{n}
\]
ובנוסף הסדרה $y_m = \frac{1}{m} \sum_{n = 1}^m x_n$ המתכנסת ל־$A \in \RR$.
נראה שגם הסדרה $x_n$ מתכנסת ל־$A$.
\begin{proof}
	נניח בלי הגבלת הכלליות ש־$C = 1, A = 0$, זאת שכן חלוקה בקבוע וחיסור בקבוע לא משפיעים על התכנסות.
	נניח בשלילה ש־$x_n$ לא מתכנסת ל־$0$, לכן קיים $\delta > 0$ כך שקיימת תת־סדרה ${\{ x_{n_k} \}}_{k = 1}^\infty \subseteq \RR$ כך ש־$x_{n_k} \ge \delta$ לכל $k$ (בלי הגבלת הכלליות).
	מהנחת השלילה על ההתכנסות לכל $k$ יש אינסוף ערכי $n > n_k$ כך ש־$x_n > \frac{\delta}{2}$, ולכן נוכל להסיק מהצפיפות שלהם ועבור $n_k$ גדול מספיק יחד עם,
	\[
		|x_{n + m} - x_n| \le \frac{m}{n}
	\]
	שבהכרח $x_n > \frac{\delta}{2}$ תמיד, כלומר אם נבחר $\epsilon = \frac{1}{n}$ עבור $n$ זה נקבל ש־$y_{n_k}$ חסומה על־ידי $\frac{\delta}{2}$ אבל $x_n > \frac{\delta}{2}$ ואף $x_{n_k} > \delta$.
	נסיק אם כך ש־$y_n \frac{\delta}{2}$ עבור $n$ גדול מספיק, וזו סתירה.
\end{proof}

\subquestion{}
נסיק שאם $f$ אינטגרבילית רימן בקטע $[-\pi, \pi]$ אז קיים $C > 0$ כך שלכל $n$,
מקדמי פורייה של $f$ מקיימים,
\[
	|a_n|, |b_n|
	\le \frac{C}{n}
\]
ואם יש $x_0 \in [-\pi, \pi]$ כך שמתקיים $K_M * f(x_0) \to A$ אז גם $D_n * f(x_0) \to A$.
\begin{proof}
	בהגדרה של אינטגרביליות רימן דרשנו חסימות, לכן נניח ש־$f$ חסומה, וכן ממטלות קודמות ושימוש בכלל נגזרת טור פורייה (תוך שימוש באינטגרציה בחלקים וגזירת האופרטורים הטריגונומטריים) נקבל חסימות כמבוקש.
	בפרט אם קיים $x_0$ כזה, אז אנו עומדים בדרישות המדייקות של סעיף א' ונסיק שאכן יש התכנסות.
\end{proof}

\question{}
\subquestion{}
נראה שאם $V$ סוף־מימדי ו־$U \subseteq V$ סגורה,
אז לכל $v \in V$ קיים $u \in U$ כך ש־$\dist(U, v) = \lVert u - v \rVert$.
\begin{proof}
	מהנתון לכל $\epsilon > 0$ קיים $u_{\epsilon} \in U$ כך ש־$\dist(U, v) \le \lVert v - u_{\epsilon} \rVert$.
	בפרט עבור $\epsilon = \frac{1}{n}$ נקבל $u_n = u_{\frac{1}{n}}$ איבר המקיים זאת, ולכן ${\{ u_n \}}_{n = 1}^\infty \subseteq U$ סדרה כך ש־$\lVert v - u_n \rVert \to \dist(U, v)$.
	נבחר $U' = U \cap \overline{B}(v, 2 \dist(U, v))$, אז $u_n \in U'$ לכמעט כל $n$, וזו קבוצה סגורה וחסומה, לכן היא קומפקטית ויש לה תת־סדרה מתכנסת $u_{n_k}$, הנקודה הגבולית שלה $u$ מקיימת $u \in U' \subseteq U$,
	ובנוסף $\lVert u - v \rVert \le \dist(U, v) + \frac{1}{n}$ לכל $n \in \NN$, לכן בפרט $\lVert u - v \rVert = \dist(U, v)$ ונובע שהמרחק אכן מתקבל.
\end{proof}

\subquestion{}
במרחב $(\RR^2, \lVert \cdot \rVert_\infty)$ נמצא דוגמה לקבוצה $U \subseteq \RR^2$ סגורה וקמורה ולאיבר $v \in \RR^2$ כך שקיימים אינסוף $u \in U$ המקיימים $\lVert u - v \rVert_\infty = \dist(U, v)$.
\begin{solution}
	נבחר את הדיסק $U = \overline{A}_1^2(0)$ ואת הנקודה $v = 0$, אז $\dist(U, v) = 1$ וכן $\lVert v - u \rVert_\infty = 1$ לכל $u \in \partial B(0, 1)$.
	נבחין כי טענה זו נכונה עבור כל נורמה, בפרט עבור $\lVert \cdot \rVert_\infty$.
\end{solution}

\end{document}
